% Physical page 2 of the concise study, which carries this sheet and nothing
% else. Every row is the generated orientation-v2 projection of
% research/chronology.toml, placed once in the generator's order, exactly as
% the reviewed Scriptural Date and Location appendix of the expansive study
% places it, with that appendix's occurrence-and-place phrases. The generated
% definitions are imported here, once, because the concise entrypoint's own
% preamble does not carry them.
% Generated. Do not edit.
% Deterministic projection of research/chronology.toml.
% schema: 5
% document: liturgy/roman-rite/1962/propers/temporal/f02-all-saints
% generated-by: tools/tpt proper-chronology annotations
\newcommand{\chronologyannotationclaim}[6]{#6}
\newcommand{\chronologyannotationreach}[3]{}
\newcommand{\chronologyannotationgroup}[3]{#3}
\newcommand{\chronologyannotation}[1]{%
  \ifcsname triptychchronologyannotation@#1\endcsname
    \csname triptychchronologyannotation@#1\endcsname
  \else
    \PackageError{triptych}{No chronology annotation for #1}{}%
  \fi}
% introit: Introit
\expandafter\def\csname triptychchronologyannotation@introit\endcsname{%
  \chronologyannotationgroup{traditional-attribution}{disputed}{\textbf{Traditional attribution} -- disputed: \chronologyannotationreach{Ps.32.1}{Ps.32}{true}\chronologyannotationclaim{israel.monarchy.david-traditional-era}{traditional-attribution}{catholic-traditional-v1}{disputed}{reigned from 1055 to 1015 B.C.}{David (reign in the usual chronology), B.C. 1055--1015}.}%
  \space \chronologyannotationgroup{historical-setting}{research-pending}{\textbf{Historical setting} -- Occasion date unresolved.}%
  \space \chronologyannotationgroup{composition}{preferred}{\textbf{Composition}: \chronologyannotationreach{Ps.32.1}{Ps}{true}\chronologyannotationclaim{critical.psalms.latest-composition-boundary}{composition}{catholic-critical-v1}{preferred}{before the Maccabean period, around 165 B.C.; no individual psalm can be dated securely}{Before c. 165 B.C.}}%
}
% lesson: Lesson
\expandafter\def\csname triptychchronologyannotation@lesson\endcsname{%
  \chronologyannotationgroup{composition}{disputed}{\textbf{Composition} -- disputed: \chronologyannotationreach{Apoc.7.10}{Apoc}{true}\chronologyannotationreach{Apoc.7.11}{Apoc}{true}\chronologyannotationreach{Apoc.7.12}{Apoc}{true}\chronologyannotationreach{Apoc.7.2}{Apoc}{true}\chronologyannotationreach{Apoc.7.3}{Apoc}{true}\chronologyannotationreach{Apoc.7.4}{Apoc}{true}\chronologyannotationreach{Apoc.7.5}{Apoc}{true}\chronologyannotationreach{Apoc.7.6}{Apoc}{true}\chronologyannotationreach{Apoc.7.7}{Apoc}{true}\chronologyannotationreach{Apoc.7.8}{Apoc}{true}\chronologyannotationreach{Apoc.7.9}{Apoc}{true}\chronologyannotationclaim{composition.apocalypse}{composition}{catholic-traditional-v1}{disputed}{the reign of Domitian ... A.D. 81-96}{A.D. 81--96}; \chronologyannotationreach{Apoc.7.10}{Apoc}{true}\chronologyannotationreach{Apoc.7.11}{Apoc}{true}\chronologyannotationreach{Apoc.7.12}{Apoc}{true}\chronologyannotationreach{Apoc.7.2}{Apoc}{true}\chronologyannotationreach{Apoc.7.3}{Apoc}{true}\chronologyannotationreach{Apoc.7.4}{Apoc}{true}\chronologyannotationreach{Apoc.7.5}{Apoc}{true}\chronologyannotationreach{Apoc.7.6}{Apoc}{true}\chronologyannotationreach{Apoc.7.7}{Apoc}{true}\chronologyannotationreach{Apoc.7.8}{Apoc}{true}\chronologyannotationreach{Apoc.7.9}{Apoc}{true}\chronologyannotationclaim{composition.apocalypse}{composition}{catholic-traditional-v1}{disputed}{from the year 90 to 100 (approximately)}{c. A.D. 90--100}.}%
}
% gradual: Gradual
\expandafter\def\csname triptychchronologyannotation@gradual\endcsname{%
  \chronologyannotationgroup{traditional-attribution}{disputed}{\textbf{Traditional attribution} -- disputed: \chronologyannotationreach{Ps.33.10}{Ps.33 Ps.39 Ps.70 Ps.85 Ps.94 Ps.97}{true}\chronologyannotationreach{Ps.33.11}{Ps.33 Ps.39 Ps.70 Ps.85 Ps.94 Ps.97}{true}\chronologyannotationclaim{israel.monarchy.david-traditional-era}{traditional-attribution}{catholic-traditional-v1}{disputed}{reigned from 1055 to 1015 B.C.}{David (reign in the usual chronology), B.C. 1055--1015}.}%
  \space \chronologyannotationgroup{superscription-setting}{preferred}{\textbf{Superscription setting}: \chronologyannotationreach{Ps.33.10}{Ps.33 Ps.55}{true}\chronologyannotationreach{Ps.33.11}{Ps.33 Ps.55}{true}\chronologyannotationclaim{israel.monarchy.david-at-geth}{superscription-setting}{catholic-traditional-v1}{preferred}{A. M. 2944, A. C. 1060}{David's flight to Achis, king of Geth, A.M. 2944}.}%
  \space \chronologyannotationgroup{composition}{preferred}{\textbf{Composition}: \chronologyannotationreach{Ps.33.10}{Ps}{true}\chronologyannotationreach{Ps.33.11}{Ps}{true}\chronologyannotationclaim{critical.psalms.latest-composition-boundary}{composition}{catholic-critical-v1}{preferred}{before the Maccabean period, around 165 B.C.; no individual psalm can be dated securely}{Before c. 165 B.C.}}%
}
% alleluia: Alleluia
\expandafter\def\csname triptychchronologyannotation@alleluia\endcsname{%
  \chronologyannotationgroup{composition}{disputed}{\textbf{Composition} -- disputed: \chronologyannotationreach{Matt.11.28}{Matt}{true}\chronologyannotationclaim{composition.gospel-of-matthew}{composition}{catholic-traditional-v1}{disputed}{about A.D. 38-45}{c. A.D. 38--45}; \chronologyannotationreach{Matt.11.28}{Matt}{true}\chronologyannotationclaim{composition.gospel-of-matthew}{composition}{catholic-traditional-v1}{disputed}{about the year 40-42}{c. A.D. 40--42}; \chronologyannotationreach{Matt.11.28}{Matt}{true}\chronologyannotationclaim{composition.gospel-of-matthew}{composition}{catholic-traditional-v1}{disputed}{the years 40-45}{A.D. 40--45}; \chronologyannotationreach{Matt.11.28}{Matt}{true}\chronologyannotationclaim{composition.gospel-of-matthew}{composition}{catholic-traditional-v1}{disputed}{about the year 60-68}{c. A.D. 60--68}; \chronologyannotationreach{Matt.11.28}{Matt}{true}\chronologyannotationclaim{composition.gospel-of-matthew}{composition}{catholic-traditional-v1}{disputed}{about the years 64-67}{c. A.D. 64--67}; \chronologyannotationreach{Matt.11.28}{Matt}{true}\chronologyannotationclaim{composition.gospel-of-matthew}{composition}{catholic-traditional-v1}{disputed}{about the year 50}{c. A.D. 50}.}%
}
% gospel: Gospel
\expandafter\def\csname triptychchronologyannotation@gospel\endcsname{%
  \chronologyannotationgroup{narrated-event}{research-pending}{\textbf{Event} -- No narrated-event date in the chronology corpus.}%
  \space \chronologyannotationgroup{composition}{disputed}{\textbf{Composition} -- disputed: \chronologyannotationreach{Matt.5.1}{Matt}{true}\chronologyannotationreach{Matt.5.10}{Matt}{true}\chronologyannotationreach{Matt.5.11}{Matt}{true}\chronologyannotationreach{Matt.5.12}{Matt}{true}\chronologyannotationreach{Matt.5.2}{Matt}{true}\chronologyannotationreach{Matt.5.3}{Matt}{true}\chronologyannotationreach{Matt.5.4}{Matt}{true}\chronologyannotationreach{Matt.5.5}{Matt}{true}\chronologyannotationreach{Matt.5.6}{Matt}{true}\chronologyannotationreach{Matt.5.7}{Matt}{true}\chronologyannotationreach{Matt.5.8}{Matt}{true}\chronologyannotationreach{Matt.5.9}{Matt}{true}\chronologyannotationclaim{composition.gospel-of-matthew}{composition}{catholic-traditional-v1}{disputed}{about A.D. 38-45}{c. A.D. 38--45}; \chronologyannotationreach{Matt.5.1}{Matt}{true}\chronologyannotationreach{Matt.5.10}{Matt}{true}\chronologyannotationreach{Matt.5.11}{Matt}{true}\chronologyannotationreach{Matt.5.12}{Matt}{true}\chronologyannotationreach{Matt.5.2}{Matt}{true}\chronologyannotationreach{Matt.5.3}{Matt}{true}\chronologyannotationreach{Matt.5.4}{Matt}{true}\chronologyannotationreach{Matt.5.5}{Matt}{true}\chronologyannotationreach{Matt.5.6}{Matt}{true}\chronologyannotationreach{Matt.5.7}{Matt}{true}\chronologyannotationreach{Matt.5.8}{Matt}{true}\chronologyannotationreach{Matt.5.9}{Matt}{true}\chronologyannotationclaim{composition.gospel-of-matthew}{composition}{catholic-traditional-v1}{disputed}{about the year 40-42}{c. A.D. 40--42}; \chronologyannotationreach{Matt.5.1}{Matt}{true}\chronologyannotationreach{Matt.5.10}{Matt}{true}\chronologyannotationreach{Matt.5.11}{Matt}{true}\chronologyannotationreach{Matt.5.12}{Matt}{true}\chronologyannotationreach{Matt.5.2}{Matt}{true}\chronologyannotationreach{Matt.5.3}{Matt}{true}\chronologyannotationreach{Matt.5.4}{Matt}{true}\chronologyannotationreach{Matt.5.5}{Matt}{true}\chronologyannotationreach{Matt.5.6}{Matt}{true}\chronologyannotationreach{Matt.5.7}{Matt}{true}\chronologyannotationreach{Matt.5.8}{Matt}{true}\chronologyannotationreach{Matt.5.9}{Matt}{true}\chronologyannotationclaim{composition.gospel-of-matthew}{composition}{catholic-traditional-v1}{disputed}{the years 40-45}{A.D. 40--45}; \chronologyannotationreach{Matt.5.1}{Matt}{true}\chronologyannotationreach{Matt.5.10}{Matt}{true}\chronologyannotationreach{Matt.5.11}{Matt}{true}\chronologyannotationreach{Matt.5.12}{Matt}{true}\chronologyannotationreach{Matt.5.2}{Matt}{true}\chronologyannotationreach{Matt.5.3}{Matt}{true}\chronologyannotationreach{Matt.5.4}{Matt}{true}\chronologyannotationreach{Matt.5.5}{Matt}{true}\chronologyannotationreach{Matt.5.6}{Matt}{true}\chronologyannotationreach{Matt.5.7}{Matt}{true}\chronologyannotationreach{Matt.5.8}{Matt}{true}\chronologyannotationreach{Matt.5.9}{Matt}{true}\chronologyannotationclaim{composition.gospel-of-matthew}{composition}{catholic-traditional-v1}{disputed}{about the year 60-68}{c. A.D. 60--68}; \chronologyannotationreach{Matt.5.1}{Matt}{true}\chronologyannotationreach{Matt.5.10}{Matt}{true}\chronologyannotationreach{Matt.5.11}{Matt}{true}\chronologyannotationreach{Matt.5.12}{Matt}{true}\chronologyannotationreach{Matt.5.2}{Matt}{true}\chronologyannotationreach{Matt.5.3}{Matt}{true}\chronologyannotationreach{Matt.5.4}{Matt}{true}\chronologyannotationreach{Matt.5.5}{Matt}{true}\chronologyannotationreach{Matt.5.6}{Matt}{true}\chronologyannotationreach{Matt.5.7}{Matt}{true}\chronologyannotationreach{Matt.5.8}{Matt}{true}\chronologyannotationreach{Matt.5.9}{Matt}{true}\chronologyannotationclaim{composition.gospel-of-matthew}{composition}{catholic-traditional-v1}{disputed}{about the years 64-67}{c. A.D. 64--67}; \chronologyannotationreach{Matt.5.1}{Matt}{true}\chronologyannotationreach{Matt.5.10}{Matt}{true}\chronologyannotationreach{Matt.5.11}{Matt}{true}\chronologyannotationreach{Matt.5.12}{Matt}{true}\chronologyannotationreach{Matt.5.2}{Matt}{true}\chronologyannotationreach{Matt.5.3}{Matt}{true}\chronologyannotationreach{Matt.5.4}{Matt}{true}\chronologyannotationreach{Matt.5.5}{Matt}{true}\chronologyannotationreach{Matt.5.6}{Matt}{true}\chronologyannotationreach{Matt.5.7}{Matt}{true}\chronologyannotationreach{Matt.5.8}{Matt}{true}\chronologyannotationreach{Matt.5.9}{Matt}{true}\chronologyannotationclaim{composition.gospel-of-matthew}{composition}{catholic-traditional-v1}{disputed}{about the year 50}{c. A.D. 50}.}%
}
% offertory: Offertory
\expandafter\def\csname triptychchronologyannotation@offertory\endcsname{%
  \chronologyannotationgroup{composition}{disputed}{\textbf{Composition} -- disputed: \chronologyannotationreach{Wis.3.1}{Wis}{true}\chronologyannotationreach{Wis.3.2}{Wis}{true}\chronologyannotationreach{Wis.3.3}{Wis}{true}\chronologyannotationclaim{composition.book-of-wisdom}{composition}{catholic-traditional-v1}{disputed}{the time of ... Ptolemy IV Philopator (221-204 B.C.)}{In or after the time of Ptolemy IV Philopator (221-204 B.C.)}; \chronologyannotationreach{Wis.3.1}{Wis}{true}\chronologyannotationreach{Wis.3.2}{Wis}{true}\chronologyannotationreach{Wis.3.3}{Wis}{true}\chronologyannotationclaim{composition.book-of-wisdom}{composition}{catholic-traditional-v1}{disputed}{the time of ... Ptolemy VII Physicon (145-117 B.C.)}{In or after the time of Ptolemy VII Physicon (145-117 B.C.)}.}%
}
% communion: Communion
\expandafter\def\csname triptychchronologyannotation@communion\endcsname{%
  \chronologyannotationgroup{composition}{disputed}{\textbf{Composition} -- disputed: \chronologyannotationreach{Matt.5.10}{Matt}{true}\chronologyannotationreach{Matt.5.8}{Matt}{true}\chronologyannotationreach{Matt.5.9}{Matt}{true}\chronologyannotationclaim{composition.gospel-of-matthew}{composition}{catholic-traditional-v1}{disputed}{about A.D. 38-45}{c. A.D. 38--45}; \chronologyannotationreach{Matt.5.10}{Matt}{true}\chronologyannotationreach{Matt.5.8}{Matt}{true}\chronologyannotationreach{Matt.5.9}{Matt}{true}\chronologyannotationclaim{composition.gospel-of-matthew}{composition}{catholic-traditional-v1}{disputed}{about the year 40-42}{c. A.D. 40--42}; \chronologyannotationreach{Matt.5.10}{Matt}{true}\chronologyannotationreach{Matt.5.8}{Matt}{true}\chronologyannotationreach{Matt.5.9}{Matt}{true}\chronologyannotationclaim{composition.gospel-of-matthew}{composition}{catholic-traditional-v1}{disputed}{the years 40-45}{A.D. 40--45}; \chronologyannotationreach{Matt.5.10}{Matt}{true}\chronologyannotationreach{Matt.5.8}{Matt}{true}\chronologyannotationreach{Matt.5.9}{Matt}{true}\chronologyannotationclaim{composition.gospel-of-matthew}{composition}{catholic-traditional-v1}{disputed}{about the year 60-68}{c. A.D. 60--68}; \chronologyannotationreach{Matt.5.10}{Matt}{true}\chronologyannotationreach{Matt.5.8}{Matt}{true}\chronologyannotationreach{Matt.5.9}{Matt}{true}\chronologyannotationclaim{composition.gospel-of-matthew}{composition}{catholic-traditional-v1}{disputed}{about the years 64-67}{c. A.D. 64--67}; \chronologyannotationreach{Matt.5.10}{Matt}{true}\chronologyannotationreach{Matt.5.8}{Matt}{true}\chronologyannotationreach{Matt.5.9}{Matt}{true}\chronologyannotationclaim{composition.gospel-of-matthew}{composition}{catholic-traditional-v1}{disputed}{about the year 50}{c. A.D. 50}.}%
}
% presentation: orientation-v2
\newcommand{\chronologyorientationrow}[2]{%
  \ifcsname triptychchronologyorientation@#1\endcsname
    \csname triptychchronologyorientation@#1\endcsname{#2}%
  \else
    \PackageError{triptych}{No chronology orientation row for #1}{}%
  \fi}
\expandafter\def\csname triptychchronologyorientation@row-e88e943a6faaa20b\endcsname#1{%
Introit: Psalms 32:1 (Vulgate numbering; Hebrew 33); Gradual: Psalms 33:10--11 (Vulgate numbering; Hebrew 34) & Traditional attribution: David's reign. #1 & Disputed: \chronologyannotationreach{Ps.32.1}{Ps.32}{true}\chronologyannotationreach{Ps.33.10}{Ps.33 Ps.39 Ps.70 Ps.85 Ps.94 Ps.97}{true}\chronologyannotationreach{Ps.33.11}{Ps.33 Ps.39 Ps.70 Ps.85 Ps.94 Ps.97}{true}\chronologyannotationclaim{israel.monarchy.david-traditional-era}{traditional-attribution}{catholic-traditional-v1}{disputed}{reigned from 1055 to 1015 B.C.}{B.C. 1055--1015}\\[.25em]%
}
\expandafter\def\csname triptychchronologyorientation@row-bbfc04633a7fe6e8\endcsname#1{%
Introit: Psalms 32:1 (Vulgate numbering; Hebrew 33) & Historical setting. #1 & Occasion date unresolved\\[.25em]%
}
\expandafter\def\csname triptychchronologyorientation@row-fb6656afd6f4b4f0\endcsname#1{%
Introit: Psalms 32:1 (Vulgate numbering; Hebrew 33); Gradual: Psalms 33:10--11 (Vulgate numbering; Hebrew 34) & Composition: Psalms. #1 & \chronologyannotationreach{Ps.32.1}{Ps}{true}\chronologyannotationreach{Ps.33.10}{Ps}{true}\chronologyannotationreach{Ps.33.11}{Ps}{true}\chronologyannotationclaim{critical.psalms.latest-composition-boundary}{composition}{catholic-critical-v1}{preferred}{before the Maccabean period, around 165 B.C.; no individual psalm can be dated securely}{Before c. 165 B.C.}\\[.25em]%
}
\expandafter\def\csname triptychchronologyorientation@row-d7a171509b4d3dc5\endcsname#1{%
Gradual: Psalms 33:10--11 (Vulgate numbering; Hebrew 34) & Superscription setting: David's flight to Achis, king of Geth. #1 & \chronologyannotationreach{Ps.33.10}{Ps.33 Ps.55}{true}\chronologyannotationreach{Ps.33.11}{Ps.33 Ps.55}{true}\chronologyannotationclaim{israel.monarchy.david-at-geth}{superscription-setting}{catholic-traditional-v1}{preferred}{A. M. 2944, A. C. 1060}{A.M. 2944 (Haydock edition reporting Ussher)}\\[.25em]%
}
\expandafter\def\csname triptychchronologyorientation@row-43dce930fad06090\endcsname#1{%
Offertory: Wisdom 3:1--3 & Composition: The Book of Wisdom. #1 & Disputed: \chronologyannotationreach{Wis.3.1}{Wis}{true}\chronologyannotationreach{Wis.3.2}{Wis}{true}\chronologyannotationreach{Wis.3.3}{Wis}{true}\chronologyannotationclaim{composition.book-of-wisdom}{composition}{catholic-traditional-v1}{disputed}{the time of ... Ptolemy IV Philopator (221-204 B.C.)}{In or after the time of Ptolemy IV Philopator (221-204 B.C.)}; \chronologyannotationreach{Wis.3.1}{Wis}{true}\chronologyannotationreach{Wis.3.2}{Wis}{true}\chronologyannotationreach{Wis.3.3}{Wis}{true}\chronologyannotationclaim{composition.book-of-wisdom}{composition}{catholic-traditional-v1}{disputed}{the time of ... Ptolemy VII Physicon (145-117 B.C.)}{In or after the time of Ptolemy VII Physicon (145-117 B.C.)}\\[.25em]%
}
\expandafter\def\csname triptychchronologyorientation@row-c1e1f88c6b954af2\endcsname#1{%
Gospel: Matthew 5:1--12 & Event. #1 & Date unresolved\\[.25em]%
}
\expandafter\def\csname triptychchronologyorientation@row-8b85083df79d6e95\endcsname#1{%
Alleluia: Matthew 11:28; Gospel: Matthew 5:1--12; Communion: Matthew 5:8--10 & Composition: The Gospel of St Matthew. #1 & Disputed: \chronologyannotationreach{Matt.11.28}{Matt}{true}\chronologyannotationreach{Matt.5.1}{Matt}{true}\chronologyannotationreach{Matt.5.10}{Matt}{true}\chronologyannotationreach{Matt.5.11}{Matt}{true}\chronologyannotationreach{Matt.5.12}{Matt}{true}\chronologyannotationreach{Matt.5.2}{Matt}{true}\chronologyannotationreach{Matt.5.3}{Matt}{true}\chronologyannotationreach{Matt.5.4}{Matt}{true}\chronologyannotationreach{Matt.5.5}{Matt}{true}\chronologyannotationreach{Matt.5.6}{Matt}{true}\chronologyannotationreach{Matt.5.7}{Matt}{true}\chronologyannotationreach{Matt.5.8}{Matt}{true}\chronologyannotationreach{Matt.5.9}{Matt}{true}\chronologyannotationclaim{composition.gospel-of-matthew}{composition}{catholic-traditional-v1}{disputed}{about A.D. 38-45}{c. A.D. 38--45}; \chronologyannotationreach{Matt.11.28}{Matt}{true}\chronologyannotationreach{Matt.5.1}{Matt}{true}\chronologyannotationreach{Matt.5.10}{Matt}{true}\chronologyannotationreach{Matt.5.11}{Matt}{true}\chronologyannotationreach{Matt.5.12}{Matt}{true}\chronologyannotationreach{Matt.5.2}{Matt}{true}\chronologyannotationreach{Matt.5.3}{Matt}{true}\chronologyannotationreach{Matt.5.4}{Matt}{true}\chronologyannotationreach{Matt.5.5}{Matt}{true}\chronologyannotationreach{Matt.5.6}{Matt}{true}\chronologyannotationreach{Matt.5.7}{Matt}{true}\chronologyannotationreach{Matt.5.8}{Matt}{true}\chronologyannotationreach{Matt.5.9}{Matt}{true}\chronologyannotationclaim{composition.gospel-of-matthew}{composition}{catholic-traditional-v1}{disputed}{about the year 40-42}{c. A.D. 40--42}; \chronologyannotationreach{Matt.11.28}{Matt}{true}\chronologyannotationreach{Matt.5.1}{Matt}{true}\chronologyannotationreach{Matt.5.10}{Matt}{true}\chronologyannotationreach{Matt.5.11}{Matt}{true}\chronologyannotationreach{Matt.5.12}{Matt}{true}\chronologyannotationreach{Matt.5.2}{Matt}{true}\chronologyannotationreach{Matt.5.3}{Matt}{true}\chronologyannotationreach{Matt.5.4}{Matt}{true}\chronologyannotationreach{Matt.5.5}{Matt}{true}\chronologyannotationreach{Matt.5.6}{Matt}{true}\chronologyannotationreach{Matt.5.7}{Matt}{true}\chronologyannotationreach{Matt.5.8}{Matt}{true}\chronologyannotationreach{Matt.5.9}{Matt}{true}\chronologyannotationclaim{composition.gospel-of-matthew}{composition}{catholic-traditional-v1}{disputed}{the years 40-45}{A.D. 40--45}; \chronologyannotationreach{Matt.11.28}{Matt}{true}\chronologyannotationreach{Matt.5.1}{Matt}{true}\chronologyannotationreach{Matt.5.10}{Matt}{true}\chronologyannotationreach{Matt.5.11}{Matt}{true}\chronologyannotationreach{Matt.5.12}{Matt}{true}\chronologyannotationreach{Matt.5.2}{Matt}{true}\chronologyannotationreach{Matt.5.3}{Matt}{true}\chronologyannotationreach{Matt.5.4}{Matt}{true}\chronologyannotationreach{Matt.5.5}{Matt}{true}\chronologyannotationreach{Matt.5.6}{Matt}{true}\chronologyannotationreach{Matt.5.7}{Matt}{true}\chronologyannotationreach{Matt.5.8}{Matt}{true}\chronologyannotationreach{Matt.5.9}{Matt}{true}\chronologyannotationclaim{composition.gospel-of-matthew}{composition}{catholic-traditional-v1}{disputed}{about the year 60-68}{c. A.D. 60--68}; \chronologyannotationreach{Matt.11.28}{Matt}{true}\chronologyannotationreach{Matt.5.1}{Matt}{true}\chronologyannotationreach{Matt.5.10}{Matt}{true}\chronologyannotationreach{Matt.5.11}{Matt}{true}\chronologyannotationreach{Matt.5.12}{Matt}{true}\chronologyannotationreach{Matt.5.2}{Matt}{true}\chronologyannotationreach{Matt.5.3}{Matt}{true}\chronologyannotationreach{Matt.5.4}{Matt}{true}\chronologyannotationreach{Matt.5.5}{Matt}{true}\chronologyannotationreach{Matt.5.6}{Matt}{true}\chronologyannotationreach{Matt.5.7}{Matt}{true}\chronologyannotationreach{Matt.5.8}{Matt}{true}\chronologyannotationreach{Matt.5.9}{Matt}{true}\chronologyannotationclaim{composition.gospel-of-matthew}{composition}{catholic-traditional-v1}{disputed}{about the years 64-67}{c. A.D. 64--67}; \chronologyannotationreach{Matt.11.28}{Matt}{true}\chronologyannotationreach{Matt.5.1}{Matt}{true}\chronologyannotationreach{Matt.5.10}{Matt}{true}\chronologyannotationreach{Matt.5.11}{Matt}{true}\chronologyannotationreach{Matt.5.12}{Matt}{true}\chronologyannotationreach{Matt.5.2}{Matt}{true}\chronologyannotationreach{Matt.5.3}{Matt}{true}\chronologyannotationreach{Matt.5.4}{Matt}{true}\chronologyannotationreach{Matt.5.5}{Matt}{true}\chronologyannotationreach{Matt.5.6}{Matt}{true}\chronologyannotationreach{Matt.5.7}{Matt}{true}\chronologyannotationreach{Matt.5.8}{Matt}{true}\chronologyannotationreach{Matt.5.9}{Matt}{true}\chronologyannotationclaim{composition.gospel-of-matthew}{composition}{catholic-traditional-v1}{disputed}{about the year 50}{c. A.D. 50}\\[.25em]%
}
\expandafter\def\csname triptychchronologyorientation@row-ca811ebc12efaaee\endcsname#1{%
Lesson: Revelation 7:2--12 & Composition: The Apocalypse of St John. #1 & Disputed: \chronologyannotationreach{Apoc.7.10}{Apoc}{true}\chronologyannotationreach{Apoc.7.11}{Apoc}{true}\chronologyannotationreach{Apoc.7.12}{Apoc}{true}\chronologyannotationreach{Apoc.7.2}{Apoc}{true}\chronologyannotationreach{Apoc.7.3}{Apoc}{true}\chronologyannotationreach{Apoc.7.4}{Apoc}{true}\chronologyannotationreach{Apoc.7.5}{Apoc}{true}\chronologyannotationreach{Apoc.7.6}{Apoc}{true}\chronologyannotationreach{Apoc.7.7}{Apoc}{true}\chronologyannotationreach{Apoc.7.8}{Apoc}{true}\chronologyannotationreach{Apoc.7.9}{Apoc}{true}\chronologyannotationclaim{composition.apocalypse}{composition}{catholic-traditional-v1}{disputed}{the reign of Domitian ... A.D. 81-96}{A.D. 81--96}; \chronologyannotationreach{Apoc.7.10}{Apoc}{true}\chronologyannotationreach{Apoc.7.11}{Apoc}{true}\chronologyannotationreach{Apoc.7.12}{Apoc}{true}\chronologyannotationreach{Apoc.7.2}{Apoc}{true}\chronologyannotationreach{Apoc.7.3}{Apoc}{true}\chronologyannotationreach{Apoc.7.4}{Apoc}{true}\chronologyannotationreach{Apoc.7.5}{Apoc}{true}\chronologyannotationreach{Apoc.7.6}{Apoc}{true}\chronologyannotationreach{Apoc.7.7}{Apoc}{true}\chronologyannotationreach{Apoc.7.8}{Apoc}{true}\chronologyannotationreach{Apoc.7.9}{Apoc}{true}\chronologyannotationclaim{composition.apocalypse}{composition}{catholic-traditional-v1}{disputed}{from the year 90 to 100 (approximately)}{c. A.D. 90--100}\\[.25em]%
}


\runninghead{Scriptural Date and Location}
\section{Scriptural Date and Location}

{\small\noindent Each passage the Mass reads or sings from Scripture, in
canonical order: what is dated, where it belongs, and the date the Scripture
chronology corpus gives. Where several texts of the Mass share one date, the
row names them all; the Introit antiphon and the orations of the feast and
of the commemorated Sunday have no biblical date.\par}
\medskip

\zlabel{triptych:concise:chronology:start}
\begin{chronologyorientationtable}
\chronologyorientationrow{row-e88e943a6faaa20b}{Two psalms of praise and of the fear of the Lord. The Clementine prints ``A psalm for David'' inside Ps 32:1, where the Hebrew has no title; Ps 33 is titled ``For David'' (33:1). No occasion is dated.}
\addlinespace
\chronologyorientationrow{row-bbfc04633a7fe6e8}{A hymn bidding the just rejoice and praise the Lord; it names no place or occasion.}
\addlinespace
\chronologyorientationrow{row-fb6656afd6f4b4f0}{The book of Psalms as a whole, not these two psalms singly.}
\addlinespace
\chronologyorientationrow{row-d7a171509b4d3dc5}{David at the court of Achis, king of Geth, among the Philistines (1~Kings 21:10--15), the occasion the psalm's title names under the name Achimelech (Ps 33:1).}
\midrule
\chronologyorientationrow{row-43dce930fad06090}{The Book of Wisdom, whose answer to the wicked (Wis 2--3) the Offertory adapts; the proposed dates are reigns of Ptolemaic kings of Egypt.}
\midrule
\chronologyorientationrow{row-c1e1f88c6b954af2}{``A mountain'' (Mt 5:1), not named, after Jesus has gone ``about all Galilee'' and the multitudes have followed him (4:23--25).}
\addlinespace
\chronologyorientationrow{row-8b85083df79d6e95}{The Gospel according to St Matthew, from which the Alleluia verse, the Gospel and the Communion are taken.}
\midrule
\chronologyorientationrow{row-ca811ebc12efaaee}{John writes the visions he saw ``in the island which is called Patmos'' (Apoc 1:9).}
\end{chronologyorientationtable}
\zlabel{triptych:concise:chronology:end}

\medskip
{\footnotesize\noindent The titles' attribution of the two psalms, the two
Ptolemaic datings proposed for Wisdom, and the source and evidence of every
date are set out passage by passage in the Scriptural Date and Location
appendix of the full study.\par}
